\clearpage
\subsection{RNF-04 — Compatibilidad multiplataforma}

\renewcommand{\arraystretch}{1.3}

\noindent
\begin{tabular}{|p{5cm}|p{10.5cm}|}
\hline
\rowcolor{secVerde}\multicolumn{2}{|l|}{\color{white}\textbf{1. INFORMACIÓN DEL CLIENTE}} \\
\hline
Nombre de la empresa  & UTEQ \\ \hline
Nombre del Contacto   & Angélica Garduño Bustamante \\ \hline
Fecha de Elaboración  & 01/06/2026 \\ \hline
\end{tabular}

\vspace{0.3cm}

\noindent
\begin{tabular}{|p{5cm}|p{10.5cm}|}
\hline
\rowcolor{secVerde}\multicolumn{2}{|l|}{\color{white}\textbf{2. INFORMACIÓN GENERAL}} \\
\hline
Nombre del Requerimiento     & RNF 04. Compatibilidad multiplataforma \\ \hline
Tipo de Requerimiento        & No funcional — Portabilidad \\ \hline
Fecha requerida de solución  & 30/07/2026 \\ \hline
Prioridad                    & Media \\ \hline
\end{tabular}

\vspace{0.3cm}

\noindent
\begin{tabular}{|p{4.5cm}|p{11cm}|}
\hline
\rowcolor{secVerde}\multicolumn{2}{|l|}{\color{white}\textbf{3. DESCRIPCIÓN DEL REQUERIMIENTO}} \\
\hline
Objetivo del Requerimiento &
Garantizar que la aplicación móvil ejecute correctamente todas sus funcionalidades, incluidas las nuevas, en Android 8.0 (API 26) o superior y en iOS 13.0 o superior. \\
\hline
Beneficio esperado &
Cobertura de la mayoría de los dispositivos en uso por el público objetivo (apicultores) sin necesidad de hardware reciente. Reduce el riesgo de exclusión por incompatibilidad. \\
\hline
Alcance &
- Declarar las versiones mínimas en \texttt{app.json}. \par
- Validar las nuevas dependencias contra las versiones mínimas. \par
- Probar en al menos un emulador de Android 8 y un simulador de iOS 13. \\
\hline
Ubicación & \url{https://github.com/MarkG777/AbejaNet} \\ \hline
Descripción de Usuarios y Responsables &
\textbf{Usuarios:} \par
\textbf{Usuario (Cliente):} ejecuta la app desde un dispositivo móvil de gama media. \par
\vspace{0.1cm}
\textbf{Responsables:} \par
\textbf{Equipo de desarrollo:} configura las versiones mínimas y ejecuta pruebas en emuladores. \par
\textbf{Expo SDK 51:} provee la base de compatibilidad multiplataforma. \\
\hline
Descripción funcional &
1. La app debe instalarse y ejecutarse en Android 8.0 o superior. \par
2. La app debe instalarse y ejecutarse en iOS 13.0 o superior. \par
3. Todas las funcionalidades nuevas (RQF 03 a RQF 05) deben verificarse en ambos sistemas operativos. \par
4. Las versiones mínimas se documentan en \texttt{app.json} y en el \texttt{README}. \\
\hline
Restricciones &
- No se garantiza compatibilidad con Android $<$ 8.0 ni iOS $<$ 13.0. \par
- Las funciones biométricas dependen del hardware del dispositivo. \\
\hline
Precedencias & Ninguna. \\ \hline
Documentación requerida & \texttt{README} actualizado con las versiones mínimas. \\ \hline
Observaciones & Esta compatibilidad ya se garantiza por defecto al utilizar Expo SDK 51. \\ \hline
\end{tabular}

\vspace{0.3cm}

\noindent
\begin{tabular}{|p{5.1cm}|p{5.1cm}|p{5.1cm}|}
\hline
\rowcolor{secVerde}\multicolumn{3}{|l|}{\color{white}\textbf{4. CONDICIONES DE PRUEBA}} \\
\hline
\rowcolor{cellGris}\textbf{Condición} & \textbf{Validación} & \textbf{Resultado} \\
\hline
La app se instala en un emulador de Android 8.0 (API 26). &
Todas las pantallas y funcionalidades deben renderizarse sin errores. &
La app inicia correctamente y permite navegar por todas las pantallas. \\
\hline
La app se instala en un simulador de iOS 13.0. &
La biometría, las notificaciones y el almacenamiento seguro deben funcionar. &
Todas las funciones críticas operan sin advertencias de compatibilidad. \\
\hline
\end{tabular}
